\section*{Bab \ref{ch:g12:brain-and-movement} --- Dari Niat Menjadi Gerak}

\begin{solution}{exo:g12:brain-and-movement:1}
Sebuah pita di puncak tiap belahan otak, di depan alur antara lobus
frontal dan lobus parietal. Petanya menugaskan tiap bagian pita itu
kepada satu bagian sisi tubuh yang berlawanan, dengan luas yang
sebanding dengan kehalusan gerak bagian tersebut.
\end{solution}

\begin{solution}{exo:g12:brain-and-movement:2}
Akson \omterm{def:g12:stretch-reflex:neuron}{neuron} \omterm{def:g12:brain-and-movement:cortex}{korteks motorik} turun lewat batang otak, tempat sebagian
besarnya menyeberang ke sisi yang lain, lalu menuruni sumsum tulang
belakang sampai aras tungkainya, tempat akson itu berakhir pada \omterm{def:g12:stretch-reflex:neuron}{neuron}
motorik (atau interneuron di sebelahnya) yang aksonnya berjalan ke otot
tungkainya.
\end{solution}

\begin{solution}{exo:g12:brain-and-movement:3}
Luasnya mengikuti ketelitian kendalinya, bukan ukurannya: tangan
membuat banyak gerakan halus yang saling bebas dan memerlukan banyak
\omterm{def:g12:stretch-reflex:neuron}{neuron} korteks; batang tubuh membuat sedikit gerakan yang kasar.
\end{solution}

\begin{solution}{exo:g12:brain-and-movement:4}
Setiap perintah kepada sebuah otot --- dari korteksnya, dari \omterm{def:g12:stretch-reflex:reflex}{lengkung
refleksnya}, dari interneuronnya --- harus lewat \omterm{def:g12:stretch-reflex:neuron}{neuron} motoriknya, yang
menjumlahkan semua masukannya lalu memicu hanya bila jumlahnya mencapai
ambang; tak ada jalan lain menuju ototnya.
\end{solution}

\begin{solution}{exo:g12:brain-and-movement:5}
Meluasnya wilayah jari para pemusik (dan para sukarelawan setelah
berhari-hari berlatih); diambil alihnya wilayah anggota badan yang
diamputasi oleh tetangganya; dan pemulihan setelah stroke yang
diperintah wilayah di sekitar lesinya.
\end{solution}

\begin{solution}{exo:g12:brain-and-movement:6}
Sekitar \qty{4}{mV}; empat masukan beruntun cepat mencapai ambang
\qty{15}{mV}; sebuah masukan penghambat yang tiba bersama masukan
perangsang meniadakannya dan tegangannya nyaris tak bergerak.
\end{solution}

\begin{solution}{exo:g12:brain-and-movement:7}
Wajah dan lengan di sebelah kiri, tungkainya kurang terkena: \omterm{def:g12:brain-and-movement:cortex}{korteks
motorik} belahan otak kanan, pada bagian bawah dan tengahnya (wilayah
wajah dan lengan), dengan wilayah tungkai di puncaknya hanya sebagian
tersentuh --- di atas penyeberangannya, sebab kekurangannya ada pada
satu sisi termasuk wajahnya.
\end{solution}

\begin{solution}{exo:g12:brain-and-movement:8}
Setelah stroke, \omterm{def:g12:stretch-reflex:reflex}{lengkung refleksnya} utuh tetapi pengaruh penghambat
korteks pada \omterm{def:g12:stretch-reflex:neuron}{neuron} motoriknya hilang, sehingga masukan kumparannya tak
menemui peredam: sentakan yang berlebihan. Saraf yang terpotong
memutuskan lengkungnya sendiri: tak ada isyarat yang mencapai ototnya,
jadi tak ada \omterm{def:g12:stretch-reflex:reflex}{refleksnya}.
\end{solution}

\begin{solution}{exo:g12:brain-and-movement:9}
Gerakan sadar kedua tungkai dan batang tubuh bawahnya hilang (traktusnya
terputus); \omterm{def:g12:stretch-reflex:reflex}{refleks} tungkainya tetap ada, bahkan berlebihan, sebab
lengkungnya terletak di bawah lesinya; serat wajahnya meninggalkan
traktusnya di batang otak, di atas lesinya, dan tak tersentuh.
\end{solution}

\begin{solution}{exo:g12:brain-and-movement:10}
Mengepalkan tinju mengirim perangsangan umum dari korteks menuruni
sumsumnya, termasuk sebagian ke \omterm{def:g12:stretch-reflex:neuron}{neuron} motorik tungkainya; tegangannya
duduk lebih dekat ke ambangnya, dan masukan kumparannya mencapai ambang
itu lebih mudah dan pada lebih banyak \omterm{def:g12:stretch-reflex:neuron}{neuron}.
\end{solution}

\begin{solution}{exo:g12:brain-and-movement:11}
Bertahun-tahun berlatih memperluas wilayah jarinya lewat \omterm{prop:g11:vision-and-brain:plasticity}{plastisitas};
memulainya sebelum usia tujuh tahun, pada masa korteksnya paling mudah
menata ulang, menghasilkan perubahan yang lebih besar daripada latihan
yang sama pada usia yang lebih tua.
\end{solution}

\begin{solution}{exo:g12:brain-and-movement:12}
Sirkuit menguat oleh pemakaian: bila lengan yang baik mengerjakan
segalanya, sirkuit yang selamat menuju lengan yang lumpuh tak pernah
diaktifkan dan tak terseleksi; memaksakan pemakaiannya membuat sirkuit
itu memicu, dan yang berhasil pun menguat. Seperti pada anak kucing,
sambungan yang dipakai merebut wilayahnya; yang tak dipakai
kehilangannya.
\end{solution}

\begin{solution}{exo:g12:brain-and-movement:13}
Wilayah sensorik tangannya, yang kehilangan masukan, dimasuki wilayah
wajah di sebelahnya, sehingga menyentuh pipinya mengaktifkan \omterm{def:g12:stretch-reflex:neuron}{neuron}
yang masih ``berarti'' tangan. Peta motoriknya terletak di sebelahnya
dan menata ulang dengan cara yang sama: bekas wilayah tangannya
diperkirakan diambil alih wilayah lengan dan wajahnya.
\end{solution}

\begin{solution}{exo:g12:brain-and-movement:14}
\omterm{def:g12:stretch-reflex:neuron}{Neuron} motoriknya rusak: lumpuh, \omterm{def:g12:stretch-reflex:reflex}{refleksnya} hilang, ototnya menyusut
(tanpa masukan saraf); korteksnya membentuk niat yang tak dapat
mencapai satu otot pun --- sebuah antarmuka yang membaca korteksnya
dapat menggerakkan sebuah prostesis atau sebuah perangsang ototnya.
Korteksnya rusak: lumpuh, \omterm{def:g12:stretch-reflex:reflex}{refleksnya} berlebihan, ototnya terpelihara
oleh \omterm{def:g12:stretch-reflex:reflex}{refleksnya}; \omterm{def:g12:stretch-reflex:neuron}{neuron} motoriknya utuh --- sebuah antarmuka dapat
merangsang sumsumnya atau \omterm{def:g12:stretch-reflex:neuron}{neuronnya} langsung, tetapi justru niatnyalah
yang hilang, dan \omterm{prop:g11:vision-and-brain:plasticity}{plastisitas} korteks di sekitarnya adalah harapan
utamanya.
\end{solution}

\begin{solution}{exo:g12:brain-and-movement:15}
Otot memang tumbuh oleh latihan, tetapi sebuah gerakan yang dipelajari
adalah perubahan pada sistem sarafnya: wilayah korteks yang
memerintahnya meluas dan sambungannya menajam (\omterm{prop:g11:vision-and-brain:plasticity}{plastisitas}), dan pola
masukan yang mencapai tiap \omterm{def:g12:stretch-reflex:neuron}{neuron} motorik --- mana yang memicu, berapa
banyak, dalam urutan apa --- diperhalus sehingga jumlah pada \omterm{def:g12:stretch-reflex:neuron}{neuron}
motoriknya menghasilkan kontraksi yang tepat pada saat yang tepat. Otot
yang sama, yang diperintah lebih baik.
\end{solution}

\begin{solution}{pb:g12:brain-and-movement:1}
\textbf{1.} Serat kortikospinalnya menyeberang di batang otak: korteks
kiri memerintah sisi kanannya.

\textbf{2.} Bagian bawah dan bagian atas seluruh pitanya: wajah dan
lengan (sisi pitanya) serta tungkai (puncaknya) --- sebuah lesi yang
melintasi \omterm{def:g12:brain-and-movement:cortex}{korteks motorik} kirinya.

\textbf{3.} Pita sensoriknya, tepat di belakang pita motoriknya
menyeberangi alur tengahnya, luput: rasanya mencapainya dengan normal.

\textbf{4.} \omterm{def:g12:stretch-reflex:reflex}{Lengkung refleks} di dalam sumsumnya utuh, dan masukan
peredam korteksnya kepada \omterm{def:g12:stretch-reflex:neuron}{neuron} motoriknya sudah hilang: masukan
kumparannya memicu \omterm{def:g12:stretch-reflex:neuron}{neuronnya} lebih mudah.

\textbf{5.} \omterm{def:g12:stretch-reflex:neuron}{Neuron} motoriknya atau sarafnya di sebelah kanan:
lengkungnya putus (tak ada \omterm{def:g12:stretch-reflex:reflex}{refleks}) dan ototnya, tanpa masukan saraf,
menyusut. Sebuah lesi di bawah \omterm{prop:g12:brain-and-movement:integration}{jalur bersama terakhirnya}, bukan di
atasnya.

\textbf{6.} $12 \times 1 + 3 \times 2 = \qty{18}{mV}$: di atas ambangnya,
jadi \omterm{def:g12:stretch-reflex:neuron}{neuronnya} memicu.

\textbf{7.} $18 - 4 = \qty{14}{mV}$: tepat di bawah ambangnya, tak
memicu --- penghambatnya menyetel langkahnya.

\textbf{8.} $0 + 6 = \qty{6}{mV}$: tak memicu; tungkainya tak bergerak.

\textbf{9.} Sebelumnya: $20 - 6 = \qty{14}{mV}$, di bawah ambangnya atau
tepat padanya --- sentakan yang sederhana. Sesudahnya: \qty{20}{mV},
jauh di atasnya --- sentakan yang kuat: peredamannyalah yang dahulu
menjaga \omterm{def:g12:stretch-reflex:reflex}{refleksnya} tetap sederhana.

\textbf{10.} $8 + 6 = \qty{14}{mV}$: kurang sedikit; dengan masukan
kumparan (atau usaha) yang sedikit lebih banyak ia memicu. Gerakannya
lemah dan ragu, memicu pada sebagian percobaan dan tidak pada yang lain.

\textbf{11.} Dari korteks yang selamat di sekitar lesinya dan dari
belahan otak yang lain, yang \omterm{def:g12:stretch-reflex:neuron}{neuronnya} membentuk sambungan baru atau
sambungan yang menguat menuju \omterm{def:g12:stretch-reflex:neuron}{neuron} motorik tungkainya.

\textbf{12.} \omterm{prop:g11:vision-and-brain:plasticity}{Plastisitas} menguatkan sirkuit yang diaktifkan: gerakan
yang dicoba memicu \omterm{def:g12:stretch-reflex:neuron}{neuron} korteks yang selamat beserta lintasannya, dan
yang menghasilkan sebuah kontraksi pun terseleksi; tungkai yang
digerakkan secara pasif hanya mengaktifkan sirkuit sensoriknya, bukan
sirkuit yang memerintah.

\textbf{13.} Sebagian kecil serat kortikospinal yang tidak menyeberang,
dan lintasan lewat batang otaknya, dapat membawa perintah dari belahan
otak kanan ke tungkai kanannya; penataan ulangnya merekrut serat itu.

\textbf{14.} Tangan menempati wilayah yang luas dan khusus yang kendali
halusnya sukar digantikan wilayah tetangganya, dan gerakannya menuntut
ketelitian jauh lebih tinggi daripada sebuah langkah; perintah tungkai
yang lebih kasar lebih mudah dibangun kembali.

\textbf{15.} Yang sama: pemulihannya bergantung pada pemakaian, dan
sirkuit yang tak dipakai akan hilang. Yang berbeda: korteks dewasa
menata ulang jauh lebih lambat dan kurang tuntas daripada korteks anak
kucing dalam masa kritisnya --- itulah sebabnya pemulihannya memakan
waktu berbulan-bulan dan tetap sebagian.

\textbf{16.} \omterm{def:g12:brain-and-movement:cortex}{Korteks motorik}, \omterm{prop:g12:brain-and-movement:pathway}{traktus kortikospinal}, penyeberangan di
batang otak, \omterm{def:g12:stretch-reflex:neuron}{neuron} motorik sumsumnya (titik temu dengan masukan
\omterm{def:g12:stretch-reflex:reflex}{refleksnya}), asetilkolina pada sambungan neuromuskularnya, dan \omterm{def:g12:respiration-fermentation:atp}{ATP} dari
\omterm{def:g10:cells-common-unit:organelle}{mitokondria} \omterm{def:g10:muscles-and-joints:muscle}{serat ototnya}.

\textbf{17.} \omterm{prop:g12:glucose-and-diabetes:hormones}{Glukagon}; dan \omterm{prop:g12:glucose-and-diabetes:hormones}{pulau pankreasnya}.

\textbf{18.} Antisipasi oleh otaknya (perintah yang disalin ke pusat
jantungnya, disertai \omterm{prop:g10:heart-lungs-effort:control}{adrenalin}) dan \omterm{def:g11:hormones-and-reproduction:hormone}{umpan balik} dari sensor ototnya
serta darahnya.

\textbf{19.} \omterm{prop:g11:vision-and-brain:plasticity}{Plastisitas}: bab tentang penglihatan dan otak, serta bab
ini; arteri yang tersumbat: bab tentang jantung dan bab tentang
kegiatan fisik serta kesehatan; imunitas: kedua bab tentang \omterm{def:g12:innate-immunity:immunity}{imunitas
bawaan} dan \omterm{def:g12:innate-immunity:immunity}{imunitas adaptif}.

\textbf{20.} \omterm{def:g12:brain-and-movement:cortex}{Korteks motorik primer} belahan otak kiri; sebuah jumlah
masukan yang mencapai \qty{15}{mV} di atas taraf istirahat dalam
beberapa milidetik; dan \omterm{prop:g11:vision-and-brain:plasticity}{plastisitas} korteksnya.
\end{solution}
